\par\medskip
\begin{center}\hypertarget{native-166-article119}{}\pdfbookmark[2]{119.}{article-119}119.\end{center}
Haud difficile esset, e nexu inter problematis nostri data atque incognitas, eius statum ad sex aequationes reducere, vel adeo ad pauciores, quum unam alteramve incognitam satis commode eliminare liceret; sed quoniam nexus ille complicatissimus est, hae aequationes maxime intractabiles evaderent; incognitarum separatio talis, ut tandem aequatio unicam tantummodo continens prodeat, generaliter loquendo\footnote{Quoties observationes ab invicem tam parum remotae sunt, ut temporum intervalla tamquam quantitates infinite parvas tractare liceat, huiusmodi separatio utique succedit, totumque problema ad solutionem aequationis algebraicae septimi octavive gradus reducitur.} pro impossibili haberi potest, multoque adeo minus problematis solutionem integram per solas operationes directas absolvere licebit.

Sed ad \emph{duarum} aequationum solutionem $X=0$, $Y=0$, in quibus duae tantum incognitae $x$, $y$ intermixtae remanserunt, utique reducere licet problema nostrum, et quidem variis modis. Haud equidem necesse est, ut $x$, $y$ sint duo ex elementis ipsis: esse poterunt quantitates qualicunque modo cum elementis connexae, si modo illis inventis elementa inde commode derivare licet. Praeterea manifesto haud opus est, ut $X$, $Y$ per functiones explicitas ipsarum $x$, $y$ exhibeantur: sufficit, si cum illis per systema aequationum ita iunctae sunt, ut a valoribus datis ipsarum $x$, $y$ ad valores respondentes ipsarum $X$, $Y$ descendere in potestate sit.
